\documentclass[11pt]{article}

\usepackage[margin=1in]{geometry}
\usepackage{booktabs}
\usepackage{enumitem}
\usepackage[hidelinks]{hyperref}

\newcommand{\code}[1]{\texttt{#1}}

\title{Gabriel Rust Rewrite: Evaluation Plan}
\author{Plan for logging, metrics, and benchmarks}
\date{\today}

\begin{document}
\maketitle

\section{Purpose}

The design note (\code{docs/dag-tokens.tex}) claims that end-to-end tokens
and pending slots keep inputs fresh without letting backlogs build up. This
plan says how to test that claim: what to log, what to measure, which flows to
run, and what to compare against. Results will be added to this document as
they come in.

\section{When}

\begin{enumerate}[nosep]
  \item \textbf{Logging first} (done). It is small, and it closes a known gap: the
        server records no errors. It also helps debugging the rest.
  \item \textbf{Then plan step 4} (done for the local server) (result events back to producers, and
        \code{Disconnected}). Stress tests make engines and clients stall and
        die; without step 4 a flow whose server dies hangs, and a benchmark
        would measure that bug. Step 4 also lets the client time the whole
        trip of an input, from capture to result, on its own clock.
  \item \textbf{Then metrics} (timestamps done) \textbf{and benchmarks}, on
        the local server. They are repeated once the remote server exists
        (plan step 6).
\end{enumerate}
A first staleness number can be taken on the local path before that: every
process on one machine shares one monotonic clock, and step 4 does not change
the input path.

\section{Logging}

Done; the design note describes it (``Logging''). Use the \code{tracing} crate, with \code{tracing-subscriber} in tests and
benchmarks. Gabriel is a library, so it only emits events; the app chooses
whether and where to print them.
\begin{itemize}[nosep]
  \item \textbf{Errors and warnings}: every error the server and client
        threads swallow today (failed deliveries, malformed outcomes, refused
        inputs, failed sends), lost connections, and engines that stopped.
  \item \textbf{Debug}: flows installed and removed, engines added and
        removed.
  \item \textbf{Nothing per input} at the default level. Per-input data goes
        to the metrics below, not to the log.
\end{itemize}

\section{Metrics}

\subsection{Timestamps on events}

Done. Timestamps travel on the events, instead of each process writing its
own file of timestamps. The events already reach the producer's app, so:
\begin{itemize}[nosep]
  \item the app sees the whole trip of each input in one place. There are no
        files to collect from several processes and join;
  \item a process killed in a stress test loses nothing, because its events
        were sent before it died;
  \item apps get the same data in production, e.g.\ to show how old a result
        is;
  \item nothing extra is written while inputs flow. A stamp costs one read of
        the clock, about 20\,ns.
\end{itemize}

All stamps come from the monotonic clock (\code{CLOCK\_MONOTONIC} on Linux),
which every process on a machine shares. Gabriel's \code{Timestamp} type reads
it; C apps use \code{clock\_gettime(CLOCK\_MONOTONIC)}.
\begin{center}
\small
\begin{tabular}{@{}lll@{}}
\toprule
Stamp & Taken by & When \\
\midrule
$t_\mathit{publish}$ & app    & just before it publishes the input \\
$t_\mathit{offer}$   & server & it made the engine's offer: the input arrived, or \\
                     &        & the engine before passed it on (\code{offered\_at}) \\
$t_\mathit{sent}$    & server & it gave the input to the engine (\code{at} of \code{Sent}) \\
$t_\mathit{end}$     & server & the offer ended: \code{at} of \code{Result}, \code{Skipped}, \\
                     &        & or \code{TimedOut} \\
                     & client & \ldots\ for an offer on a crossing that ended on the client \\
$t_\mathit{received}$ & client & it received an event from the remote server \\
                     &        & (\code{received\_at}) \\
\bottomrule
\end{tabular}
\end{center}
\code{Producer::publish\_input} returns the input's id, so the app can match
its $t_\mathit{publish}$ to the events that carry the id. A server stamps a
result and the offers it makes with the same time, so for the next engine
$t_\mathit{offer}$ equals the result's $t_\mathit{end}$.

Not stamped yet:
\begin{itemize}[nosep]
  \item \textbf{Capture}: only the app knows when it captured a frame. It
        records that itself.
  \item \textbf{On the client}: when an input got its token, and when it was
        sent on a crossing. For the first engine, $t_\mathit{offer} -
        t_\mathit{publish}$ covers both, plus the trip to the server.
  \item \textbf{In the engine}: when it started on an input. $t_\mathit{end} -
        t_\mathit{sent}$ includes the time an input waits at the engine,
        which happens only when the server gives an engine more than one input
        at a time.
  \item \textbf{Remote server}: its stamps are on its own clock. Differences
        between two of its stamps stay valid; $t_\mathit{offer} -
        t_\mathit{publish}$ does not, across machines. The client's
        $t_\mathit{received}$ gives the AoI at the app on the device's clock
        (``Timestamps across machines'' in the design note). So far both
        servers run on one machine, so every stamp compares directly.
\end{itemize}

\subsection{Recording}

Done (\code{bench/}). The benchmark binary subscribes to every producer and keeps one row per
event: the input id, flow, engine, kind, \code{forward}, $t_\mathit{offer}$,
and the event's own time. It also keeps one row per published input with
$t_\mathit{publish}$, so inputs that never got a token show up as drops. Rows
stay in memory and are written when the run ends.

\paragraph{Format: CSV.} Written with the \code{csv} crate and \code{serde},
in long form (one row per event, not one column per stage), which is what
pandas and polars work with best. The other options:
\begin{itemize}[nosep]
  \item \textbf{Parquet} has types and compression, and loads faster. But the
        Rust crates (\code{arrow}, \code{parquet}) are large, and our runs are
        small. 30 frames a second through three engines for five minutes is
        about 50{,}000 rows, a few MB. A stress run at 1000 inputs a second
        is a few million rows, which polars still reads from CSV in seconds.
        If files grow past a few hundred MB, switch, or convert to Parquet in
        Python once.
  \item \textbf{JSON lines} are larger and slower than CSV, and our rows are
        flat, so they gain nothing.
  \item \textbf{Prometheus-style metrics} (the \code{metrics} crate) keep
        only counters and histograms. They suit a live dashboard, but every
        new question would need a new run.
\end{itemize}
Next to the rows, each run writes a small TOML file with its settings: the
flow spec, the knob values, and the git commit. Rows are never separated from
what produced them.

\subsection{What to compute}

The objective is the one in the design note (``What to optimize''): keep
waiting small while keeping the slowest engine busy. The measures below use
queueing terms; the design note defines them. \code{analyze.py} computes the
ones marked \textbf{(done)}, all from the recorded events, without the first
and last second of each run: at the start engines are warming up, and at the
end no newer input comes to replace a waiting one.
\begin{description}[style=nextline, leftmargin=1.5em]
  \item[Age of information (AoI) (done)] Per engine: at each moment, the age
        of the input behind its newest result. Its average over time, and its
        peaks just before each new result. This is the one number that
        combines both sides: it grows with waiting, and it grows when results
        are rare. Every comparison starts with it.
  \item[Waiting (done)] How old an input is when an engine gets it,
        $t_\mathit{sent} - t_\mathit{publish}$, and the \emph{excess}: this
        age minus the time the input spent in upstream engines on its path,
        $\sum (t_\mathit{end} - t_\mathit{sent})$. The excess is all the
        waiting Gabriel added. It is split by place: before the server (for a
        token, and on the crossing), $t_\mathit{offer} - t_\mathit{publish}$
        at the first engine; on an uplink, from the result of the engine
        before to the offer on the remote server (the client's wait and the
        trip out); and at each engine's edge on the server,
        $t_\mathit{sent} - t_\mathit{offer}$.
  \item[Return trip (done)] For a remote engine, $t_\mathit{received} -
        t_\mathit{end}$: how long its result took to reach the client. A
        remote result's AoI counts from $t_\mathit{received}$.
  \item[Load (done)] Per engine: the arrival rate $\lambda$ of offers, the
        throughput $X$ in results per second, the service time $S$ with its
        mean and coefficient of variation, the offered load $\rho = \lambda
        \cdot E[S]$, and the utilization $U$. A high $\rho$ with $U$ near 1
        means the engine is always busy and most candidates are skipped. A
        low $\rho$ means the engine is starved, or the flow is light.
  \item[Idle time (done)] Per engine, the gap from a result to its next
        input. This is the throughput side of the trade-off. Later: how much
        of it passed while an input waited somewhere upstream, which is the
        idle time a better limit could remove.
  \item[Waste (partly done)] The share of offers skipped in pending slots,
        and the offers made per result. Later: the compute spent upstream on
        inputs that a downstream engine skipped, and over the network, the
        bytes sent for inputs that were skipped.
  \item[Tokens in use (done)] Per producer, by Little's law: the average
        number of its inputs in flight, and the most at once, against its
        token count. An average far below the count means tokens did not
        limit the run.
  \item[Drops and why] The share of captured frames that never get a result,
        split by reason: replaced while waiting for a token, \code{Skipped}
        in a pending slot, \code{TimedOut}, \code{Disconnected}. A system can
        look fresh by dropping nearly everything, so freshness is always
        reported together with the share of frames processed.
  \item[Jitter] The spread of the time between results. Steady timing
        matters for apps such as AR.
  \item[Timeouts] How often offers time out as load grows. This tells
        whether the deadlines the compiler derives are realistic.
  \item[Fairness] When several flows or clients share an engine, the share
        of the engine each one gets.
  \item[Overhead] The latency Gabriel adds per hop (0.2\,ms at the median from
        producer to engine so far), and the CPU it uses, also when idle: the
        threads wake every 10\,ms, and control calls poll every 1\,ms.
\end{description}

\section{Benchmarks}

\subsection{Flows}

Engines in benchmarks are synthetic: they sleep or spin for a time drawn from
a chosen distribution, and report a chosen outcome. That makes every run
repeatable.
\begin{itemize}[nosep]
  \item a slow engine at the end of a deep chain;
  \item wide fan-out, with branches of very different speeds;
  \item several clients and flows sharing one engine;
  \item bursty producers, and producers faster than any engine;
  \item engine times with a heavy tail;
  \item an engine killed and restarted during a run, and the local server
        killed and restarted (needs plan step 4);
  \item with the remote server: a link that slows down and drops (possible
        now, see ``Tools'').
\end{itemize}
Each flow is run while sweeping one knob: the token count, the producer rate,
the engine times, or the number of clients.

\subsection{Baselines}

The design is only interesting next to simpler ones, run on the same flows:
\begin{itemize}[nosep]
  \item a plain FIFO queue in front of each engine (no drops, so latency
        grows without bound);
  \item keep only the newest input at each engine, with no tokens;
  \item the original Gabriel (one token per engine), where a flow can be
        expressed in it. \textbf{Done} for one engine on one machine
        (``Comparison with the original Gabriel'' below).
\end{itemize}

\subsection{Tools}

\begin{itemize}[nosep]
  \item \textbf{Done:} the benchmark binary, \code{gabriel-bench} in
        \code{bench/}. A run's TOML file holds the flow spec, each producer's
        rate (periodic or Poisson) and input size, and each engine's service
        time (constant, uniform, exponential, or log-normal), its pass
        probability, and its result size. It runs the server, the engines,
        the client, the producers, and a recorder per producer on threads of
        one process, all talking over iceoryx2 as separate processes would.
        Examples are in \code{bench/runs/}:
        \code{cargo run --release -p gabriel-bench -- bench/runs/chain.toml}.
  \item \textbf{Done:} the remote server. An engine with
        \code{location = "remote"} runs on a second server, which the client
        reaches over QUIC on loopback. A \code{[remote]} section sets the
        crossing capacity, and a \code{[network]} section the delay each way,
        jitter, loss, and rate, which the bench applies with \code{tc netem}
        on loopback. That would slow every program on the machine, so the
        bench refuses unless the run has a network namespace of its own,
        which needs no root:
        \code{unshare --user --map-root-user --net target/release/gabriel-bench
        bench/runs/remote.toml}. Only QUIC crosses loopback; iceoryx2 uses
        shared memory. Both directions get the same settings.
  \item \textbf{Done:} \code{bench/analyze.py} prints the measures marked
        done above: per producer, tokens in use; per engine, load, waiting,
        idle time, and AoI. It runs with \code{uv run}, which installs
        pandas.
  \item \textbf{Done:} \code{bench/compare/}, which compares this Gabriel
        with the original one (\code{gabriel-server} and
        \code{gabriel-client} 5.2 from PyPI) on one machine. Each system runs
        as three processes: a server, an engine that answers at once, and an
        app that sends a raw frame and waits for its result before sending
        the next. It reports the round trip and the CPU of each process:
        \code{python compare.py}.
  \item \code{hdrhistogram} for latency histograms during a run.
  \item \code{criterion} for micro-benchmarks of single parts, such as the
        flow compiler and postcard encoding.
  \item Python (pandas, matplotlib) for the analysis and the figures.
\end{itemize}

\section{Results}

\subsection{Token sweep}

How many tokens should a producer have? Each flow below was run with
several token counts on the local server: 8\,s per run without the first and
last second, one seed, one laptop. These are early numbers, not final
measurements, but the effects are large and consistent. Times are in ms.

\paragraph{Overload:} 200 inputs a second (Poisson) into one engine of
30\,ms on average with a heavy tail (log-normal, standard deviation 60\,ms;
measured $S$: 40\,ms mean, coefficient of variation 1.9).
\begin{center}
\small
\begin{tabular}{@{}rrrrrrrr@{}}
\toprule
Tokens & results/s & $U$ & $\rho$ & skipped & excess p50 & excess p99 & AoI \\
\midrule
1  & 26.2 & 0.97 & 1.1 & 0    & 3.0 & 17  & 105 \\
2  & 26.2 & 1.00 & 1.1 & 0    & 18  & 398 & 147 \\
3  & 26.2 & 0.98 & 7.9 & 0.87 & 2.1 & 20  & 103 \\
10 & 26.2 & 0.98 & 7.8 & 0.87 & 2.5 & 17  & 104 \\
\bottomrule
\end{tabular}
\end{center}

\paragraph{Balanced chain:} three engines of 20\,ms each, inputs at 100 a
second. At the last engine:
\begin{center}
\small
\begin{tabular}{@{}rrrrrr@{}}
\toprule
Tokens & results/s & $U$ & idle p50 & excess p50 & AoI \\
\midrule
1  & 16.2 & 0.33 & 41  & 5.3 & 97 \\
2  & 32.6 & 0.66 & 3.1 & 5.3 & 83 \\
3  & 48.9 & 0.99 & 0.0 & 5.6 & 77 \\
7  & 49.1 & 1.00 & 0.0 & 5.0 & 76 \\
10 & 48.9 & 1.00 & 0.0 & 5.5 & 77 \\
\bottomrule
\end{tabular}
\end{center}

\paragraph{Chain} (\code{chain.toml}): a camera at 30 frames a second, a
2\,ms filter that passes half, a 40\,ms detector with a heavy tail, and a
5\,ms tracker. At the detector:
\begin{center}
\small
\begin{tabular}{@{}rrrrrr@{}}
\toprule
Tokens & results/s & $\rho$ & excess p99 & in flight (mean) & AoI \\
\midrule
1  & 12.7 & 0.53 & 32 & 0.66 & 100 \\
2  & 14.2 & 0.60 & 91 & 0.92 & 103 \\
3  & 14.1 & 0.64 & 54 & 0.88 & 101 \\
7  & 14.1 & 0.64 & 54 & 0.88 & 101 \\
10 & 14.1 & 0.64 & 54 & 0.88 & 101 \\
\bottomrule
\end{tabular}
\end{center}

\paragraph{Fan-out} (\code{fanout.toml}): the same for 1 to 10 tokens. The
fast engine ($U = 0.16$) gets every frame, with an excess under 2\,ms at p99
and an AoI of 22\,ms. The slow one is always busy ($\rho = 5.5$, $U = 0.99$)
and skips 82\% of its offers. The token returns as soon as the fast path
completes, so the slow engine never holds it.

\paragraph{What the sweep shows,} in the terms of the design note's
trade-off: waiting, idle time, and waste.
\begin{enumerate}[nosep]
  \item \textbf{Below the knee, tokens cost idle time.} With one token, only
        one input is in the balanced chain at a time, so each engine idles
        two thirds of the time. Excess stays the same, but results come a
        third as often, and the AoI is 97 instead of 77.
  \item \textbf{Exactly enough tokens to fill the pending slots freezes
        them, which costs waiting.} With 2 tokens and one engine, one input
        is in work and one in the pending slot, and no token is left to bring
        a newer input. The waiting input waits for the rest of the engine's
        time, and the heavy tail of that time becomes the tail of the excess:
        398 at p99. Utilization gains 3\%, and the AoI rises by 40\%.
  \item \textbf{One more token unfreezes them, at the cost of waste.} Newer
        inputs keep reaching the slot and replacing the input there:
        $\rho$ jumps from 1 to 8, and 87\% of offers are skipped. On one
        machine a skipped offer costs a copy.
  \item \textbf{For one engine on the device, one token is as good.} It
        gives the same AoI as 3 tokens with no waste: fetching the next input
        when the engine is free costs only 0.3\,ms of idle time at the
        median. One token fails only when a path has several engines.
  \item \textbf{Above the bound, tokens do nothing.} With 7 or 10 tokens,
        0.9 inputs were in flight on the chain on average, and every measure
        matched 3 tokens.
  \item \textbf{Some waiting is ordinary queueing.} The chain's detector is
        not overloaded: $\rho = 0.64$. Its inputs arrive at random, as the
        filter passes half the frames, and its service time varies, so an
        input sometimes arrives while the detector is busy and waits, 54 at
        p99. No token count changes this, and neither does the AoI, which
        stays near 100 for every count. Only less load or less variance would
        lower it.
\end{enumerate}

\paragraph{Conclusion.} Tokens are not the tool that keeps inputs fresh. The
design note (``Are tokens the right tool?'') draws the consequences: tokens
become a cap above the knee, set by the rule in ``Choosing $C$''; on the
device the per-hop limits do the work; and over the network the crossing
capacity is the control to tune, by an adaptive capacity or by predictive
pull, to be compared on emulated links.

\subsection{First runs with the remote server}

\code{bench/runs/remote.toml}: a camera at 30\,fps feeds a filter on the
device (5\,ms, passes a third), which feeds a detector on the remote server
(exponential, mean 50\,ms), over a link with 20\,ms each way and 2\,ms of
jitter, and a crossing capacity of 3. Inputs are 64\,KB. Times in ms,
without the first and last second.

\paragraph{netem's jitter reorders packets, and QUIC slows down for it.} The
first runs showed inputs waiting hundreds of milliseconds on the uplink. The
cause: with jitter and no rate, netem gives each packet its own delay, so
packets overtake each other. QUIC takes a packet that arrives three behind
for lost, and cuts its sending rate. The larger the input, the more packets
it takes, and the worse it gets. Uplink time, from the filter's result to the
offer at the remote server:
\begin{center}
\small
\begin{tabular}{@{}lrr@{}}
\toprule
Run & p50 & p90 \\
\midrule
1\,KB inputs, jitter reorders     & 20  & 22 \\
16\,KB inputs, jitter reorders    & 38  & 99 \\
64\,KB inputs, jitter reorders    & 321 & 466 \\
64\,KB inputs, no jitter          & 28  & 36 \\
64\,KB inputs, jitter in order    & 30  & 36 \\
\bottomrule
\end{tabular}
\end{center}
netem keeps packets in order when it also has a rate, so the bench always
gives it one, far above what a run sends unless the run sets its own. Real
links reorder much less than this, but a lossy wireless link would hurt QUIC
the same way, through real loss. It is worth measuring with \code{loss}
next.

\paragraph{The example, with packets in order} (20\,s):
\begin{center}
\small
\begin{tabular}{@{}lrrr@{}}
\toprule
Detector & p50 & p90 & p99 \\
\midrule
uplink (client's wait and the trip out) & 28 & 33 & 46 \\
wait at its edge                        & 0  & 42 & 135 \\
idle between inputs                     & 39 & 184 & 329 \\
return trip of the result               & 20 & 22 & 22 \\
\bottomrule
\end{tabular}
\end{center}
The uplink is the 20\,ms delay plus about 8\,ms to send 64\,KB, and the
result comes back in the 20\,ms delay. The detector's load is $\rho = 0.52$
and its utilization 0.47. It idles mostly because inputs are rare (the filter
passes a third of 30 frames a second), and sometimes waits because its
service time varies. Its mean AoI at the app is 210, against 22 for the
filter on the device. A crossing capacity sweep, a slower link, and loss come
next.

\paragraph{Next in the analysis.} Compare each input's age with that of the
newest input that had reached its slot, to separate waiting Gabriel could
have avoided from waiting it could not. Measure the compute wasted upstream,
which matters once engines share a GPU.

\subsection{Comparison with the original Gabriel}

How much does each Gabriel itself cost, with everything on one machine? An
app sends a raw RGB frame of noise, one engine answers at once with 256
bytes, and the app sends the next frame when the answer arrives: one token,
one frame in flight. The engine does no work, so the round trip is all
framework. 300 frames after 30 of warm-up, one run per system, on an Intel
Core Ultra 7 265K. The original Gabriel runs in four setups:
\begin{description}[nosep, font=\normalfont\itshape]
  \item[WebSocket] The app reaches the server by WebSocket, and the engine
        reaches it by gRPC, both over TCP. This is its default setup.
  \item[WebSocket, uncompressed] The same, with the WebSocket's compression
        turned off.
  \item[ZeroMQ, TCP] The app uses ZeroMQ over TCP; the engine as above.
  \item[ZeroMQ, IPC] The app uses ZeroMQ, and both the app and the engine
        reach the server over Unix sockets.
\end{description}

\paragraph{Round trip,} median (p99) in ms:
\begin{center}
\small
\begin{tabular}{@{}lrrr@{}}
\toprule
System & 720p (2.8\,MB) & 1080p (6.2\,MB) & 4K (24.9\,MB) \\
\midrule
This Gabriel              & \textbf{0.76} (1.3) & \textbf{2.2} (3.3)
                          & \textbf{6.2} (7.9) \\
Original, ZeroMQ, IPC     & 7.6 (9.8)   & 16.2 (21.1)  & 102.6 (113.6) \\
Original, ZeroMQ, TCP     & 7.4 (10.0)  & 15.9 (19.4)  & 107.7 (115.2) \\
Original, WebSocket, uncompressed
                          & 13.3 (18.4) & 31.1 (37.2)  & 80.6 (122.4) \\
Original, WebSocket       & 67.0 (73.7) & 148.0 (157.6) & 605.3 (624.5) \\
\bottomrule
\end{tabular}
\end{center}

\paragraph{CPU,} in percent of one core, server / engine / app:
\begin{center}
\small
\begin{tabular}{@{}lrrr@{}}
\toprule
System & 720p & 1080p & 4K \\
\midrule
This Gabriel              & 38 / 8 / 57  & 32 / 6 / 62  & 39 / 2 / 58 \\
Original, ZeroMQ, IPC     & 37 / 42 / 25 & 38 / 42 / 24 & 48 / 22 / 38 \\
Original, ZeroMQ, TCP     & 35 / 42 / 27 & 37 / 44 / 27 & 50 / 21 / 36 \\
Original, WebSocket, uncompressed
                          & 45 / 23 / 37 & 48 / 23 / 35 & 40 / 24 / 43 \\
Original, WebSocket       & 6 / 7 / 88   & 6 / 6 / 90   & 10 / 4 / 88 \\
\bottomrule
\end{tabular}
\end{center}
The processes of this Gabriel do 7 to 14 times more round trips a second, so
its CPU per frame is far lower than the percentages suggest.

\paragraph{What the comparison shows.}
\begin{enumerate}[nosep]
  \item \textbf{This Gabriel is 7 to 13 times faster} than the original's
        best setup at every frame size.
  \item \textbf{The original's default setup compresses every frame.} The
        \code{websockets} library compresses each message unless told
        otherwise, and the original client has no option to turn it off. The
        app then spends 88\% of a core compressing, here on noise, which does
        not shrink. Real frames would shrink a little, for the same CPU.
  \item \textbf{The original cannot carry a raw frame over 4\,MB.} Its engine
        fails on the first 1080p frame with gRPC's default limit, which
        \code{EngineRunner} never raises:
        \begin{center}
        \code{RESOURCE\_EXHAUSTED}\\
        \code{Received message larger than max (6220814 vs. 4194304)}
        \end{center}
        The comparison raises the limit in the engine process.
  \item \textbf{Unix sockets do not help the original.} Its time goes to
        protobuf, gRPC, and asyncio, not to the socket.
  \item \textbf{Most of this Gabriel's time is spent copying the frame.}
        The server uses 0.30, 0.73, and 2.5\,ms of CPU per frame at 720p,
        1080p, and 4K: two copies of the frame at about 18\,GB/s, the speed
        of \code{memcpy}. When idle it uses about 1.5\% of a core. The four
        copies of the local path, plus the app's own, add up to most of the
        6.2\,ms at 4K (``Copies'' in the design note).
\end{enumerate}

\paragraph{Before quoting these numbers.}
\begin{itemize}[nosep]
  \item One run per system, and runs vary: the original over ZeroMQ and TCP
        took 15.9\,ms at 1080p here, and 23.3\,ms in a shorter run. Repeat
        each run and give confidence intervals.
  \item This Gabriel's engine never reads its frame, which stays in shared
        memory, while the original's engine has already received and decoded
        it. A real engine reads its frame once either way. Next: an engine
        that reads every byte, and YOLO.
  \item The original's gRPC transport for apps is not measured yet.
  \item This is one machine. The comparison that matters for Gabriel is over
        a network, with QUIC against the original's transports on emulated
        links.
\end{itemize}

\end{document}
